% ED26-BM-RAT2021-5-6: Woods's exact 2021 chapter, registered pp. 367--378.
\chapter{Reasoning, Arguments in Use, and Inference Standards}
\label{sec:p18-rationality-5-6}

Knowing that a conclusion follows, recognizing that relationship, and
actually drawing the conclusion are different matters. John Woods's
2021 chapter develops that distinction through a history and critique
of argumentation theory since the late 1950s. It examines Toulmin,
Hamblin, the Woods--Walton approach, informal logic, pragma-dialectics
and later schemes/computational developments
\cite[pp.~367--374]{Woods2021ReasoningArgumentation}.
The account is deliberately selective: its opening excludes a fuller
treatment of speech communication, rhetoric and critical thinking.
It is not a claim to cover every argumentation tradition.

\section{Consequence, recognition, and drawing}

\begin{hbdefinition}{def:p18-r56-three-consequence-roles}{Three questions about a consequence}
Fix a language and a specified consequence relation. Write
\(\Gamma\models\varphi\) for \emph{consequence-having}: the
conclusion follows from the premises under the given semantics.
This is a relation between premises and conclusion, without naming
a reasoner.

\emph{Consequence-spotting} concerns an agent's recognition of that
relationship. \emph{Consequence-drawing} concerns an agent's actual
inferential step or response. Write
\(\mathsf{Spot}(a,\Gamma,\varphi)\) and
\(\mathsf{Draw}(a,\Gamma,\varphi)\) for the corresponding agent-relative
questions; these symbols are not a new complete theory of cognition.
An agent state, goal and context must be supplied to interpret them;
their suppression in this short notation does not make them irrelevant.
Evidence that an agent reports a consequence is not automatically
evidence of correct recognition. Even correct recognition need not
settle whether drawing it now serves the agent's purpose.

Woods distinguishes an \emph{argument as such}, a finite
premise--conclusion presentation assessed by its consequence relation,
from an \emph{argument in use}, engaged by agents with purposes and
contexts. His restricted as-such description requires the terminal
member to follow from the others; it must not be mistaken for a stipulation that every passage
ordinarily called an argument is valid.
\cite[pp.~367--368]{Woods2021ReasoningArgumentation}
\index{consequence!having, spotting, drawing}\index{argument!in use}
\end{hbdefinition}

An argument can seek agreement about truth, probability, plausibility,
persuasiveness or what would be prudent to do. Those targets are not
one interchangeable success condition. Woods treats the everyday term
\emph{argument} as open textured: editorials, legal submissions,
negotiations, visual messages and quarrels share only some features.
Selecting a manageable class is legitimate, but does not establish
that its theory is canonical for every other class. His rejection of
a single exhaustive ordinary-language definition is a conceptual
position, not a proof that explicit finite formal systems cannot be defined.

The chapter's historical framing connects Aristotle's logic with fallacy
analysis and contrasts that with the later foundations-of-arithmetic
project. Woods gives nonmonotonic, relevant and paraconsistent labels
to his reconstruction of early syllogistic, with an intuitionistic
approximation as well. These source-specific labels are not properties
that can be transferred without definitions to the modern categorical
fragments at \cref{sec:p02-syllogistic-inference}. In particular their
ordinary set semantics uses the usual model-class consequence relation.
Woods's note 1 calls a property semidecidable when a finite procedure
recognizes positive instances; a decision procedure must also terminate
on negative instances. The modern finite saturation result has that
stronger stated scope
(see \cref{def:p02-syllogistic-semantics,thm:p02-syllogistic-saturation}).
It is not silently identified with every ancient presentation or with
the chapter's cited reconstruction by Corcoran.
Woods also says that decision proofs are generally not sought or available
in the broad approaches he surveys. An open-textured subject alone
defines no decision problem: its input encoding and verdict must first
be specified. That absence is not an undecidability theorem for every
particular formal argumentation system.

\section{Reasoning is not simply an argument spoken privately}

The chapter opposes two pictures. On its One picture, arguing is reasoning
aloud and reasoning is private self-argument. Its Many picture allows
reasoning that is neither addressed to another person nor internally
organized as a case offered for acceptance. A search for conclusions
from available information can differ from a search for premises in
support of an already desired conclusion. Pinto's account of an argument
as an invitation to draw an inference is one relevant view, not a
proof that every inference needs an invitation.

Woods invokes prelinguistic inference and the prevalence of uninvited
inferences in support of this distinction. The selected chapter supplies
no developmental task, recruitment, sample, design, response measure,
effect estimate or uncertainty for those claims. They remain conceptual
motivations and source assertions, not independently characterized
developmental findings or measured population frequencies. The formal
distinction between a relation and a policy does not depend on inventing
such data \cite[p.~368]{Woods2021ReasoningArgumentation}.

\begin{hbtheorem}{thm:p18-r56-consequence-policy}{proved=annex:p18.r56.consequence-policy}{A consequence relation does not select a repair policy}
For any fixed class of interpretations with ordinary all-premises
satisfaction, its semantic consequence relation is monotone:
if \(\Gamma\models\varphi\) and \(\Gamma\subseteq\Delta\),
then \(\Delta\models\varphi\).

This does not determine which beliefs an agent should retain after
discovering inconsistency. In classical propositional logic the set
\[
 K=\{p,\ p\to q,\ \neg q\}
\]
has three different maximal consistent subsets, obtained by deleting
one member. They respectively retain \(\neg q\), retain both \(p\)
and \(\neg q\), or entail \(q\). The consequence relation alone
does not choose among those repairs; a preference, reliability judgment,
goal or other revision policy is additional input.
\end{hbtheorem}

The result separates static validity from a task such as revising a belief,
seeking useful information or deciding where to spend attention. It does
not say that validity is irrelevant to reasoning, or that a sound
deductive rule becomes logically unsound because an agent has another
goal. The additional policy must itself be justified at its claimed
normative or empirical standing.

The source also challenges idealizations that require full consequence
closure, full possession of relevant information or indefinitely divisible
subjective utilities. It asks why descriptions of ideal agents should
bind actual agents who cannot instantiate them, while allowing
approximable standards implicit in practice. Whether an approximation
is meaningful depends on a specified task and representation, not on
calling a model ideal.

\begin{hbtheorem}{thm:p18-r56-explicit-closure}{proved=annex:p18.r56.explicit-closure}{An explicit-storage limit, not an impossibility of implicit representation}
In a propositional language with distinct variables
\(p,q_0,q_1,\ldots\), the consequence closure of \(\{p\}\)
contains every syntactically distinct \(p\lor q_i\).
A store of at most \(m<\omega\) explicit formula records cannot
contain that full closure.
This does not prevent representing the premises or a generating
schema finitely, answering selected consequence queries, or using
a bounded approximation. It is not a measured claim about human
memory or a proof that no finite agent can approximate any rational norm.
\end{hbtheorem}

\section{Toulmin's turn to cases made for an audience}

In Woods's account, introductory formal logic mapped selected
natural-language arguments to formal structures and used their properties
to assess the original. Toulmin's 1958 departure emphasized arguments
in use: making a case for another person's consideration. The chapter
connects that change with speech communication and, in note 3, with
Perelman and Olbrechts-Tyteca's contemporaneous rhetorical work. It
does not make one mapping or one purpose universal.

\begin{hbdefinition}{def:p18-r56-toulmin-roles}{The six roles in Woods's Toulmin discussion}
A \emph{claim} \(C\) is the conclusion being defended.
\emph{Grounds} \(G\) are offered data, facts or evidence.
A \emph{warrant} \(W\) licenses the connection from grounds to claim.
\emph{Backing} \(B\) supports a warrant when it is questioned.
A \emph{rebuttal} \(R\) marks a relevant exception or qualification;
Woods describes it as a hedging device affecting the security of the
support. A \emph{qualifier} \(Q\) expresses the limited force or
confidence of the offered conclusion.
The core of the described model comprises claim, grounds and warrant;
the other roles appear as the context requires. This is the selected
source's role analysis, not a definition of every ordinary argument.
\cite[p.~369]{Woods2021ReasoningArgumentation}
\end{hbdefinition}

Woods's illustration gives Nigel's British citizenship as claim,
birth in Bermuda as grounds, and a stated citizenship generalization
as warrant. It then supplies a story about Nigel's legal training and
knowledge, followed by an exception about betrayal and spying; a
confidence qualifier can be attached as the context requires. These
are the source's illustrative premises, not current citizenship law.
Its printed discussion marks the legal-training sentence \(R\).
Its content is support concerning the warrant, whereas the later
exception has a different role. This treatment distinguishes those
contents rather than silently passing off a repaired diagram as
Toulmin's original 1958 example.

The substantive question is why a warrant does not simply become another
premise. Woods acknowledges that the grounds alone require an appropriate
background connection, but argues that not everything needed for
successful reasoning is necessarily suitable for assertion as a premise.
His examples include tacit background information and memory; note 8
asks whether such resources must be decomposable into separate
propositionally expressible units. This is a limit on a universal
premise-list model, not a ban on making a particular expressible
meaning assumption explicit.

\begin{hbtheorem}{thm:p18-r56-lexical-representation}{proved=annex:p18.r56.lexical-representation}{The red-shirt inference has explicit formal representations}
Let \(t\) name an object and let \(\mathsf{Red},\mathsf{Colored}\)
be unary predicates. With no constraint relating their extensions,
\(\mathsf{Red}(t)\not\models\mathsf{Colored}(t)\).
With the meaning postulate
\(\forall x(\mathsf{Red}(x)\to\mathsf{Colored}(x))\),
the inference is valid. Equivalently it is valid in the restricted
model class in which the extension of \(\mathsf{Red}\) is included
in the extension of \(\mathsf{Colored}\).

If redness is explicitly represented as being colored with a red hue,
the inference has the universally valid form
\[
 \mathsf{Colored}(t)\land\mathsf{RedHue}(t)
       \models\mathsf{Colored}(t).
\]
These constructions refute an unrestricted claim that the inference
cannot be represented in any logistic system. They do not make one
semantic decomposition canonical for every context or establish that
all tacit background information has a finite premise representation.
\end{hbtheorem}

The source's note 7 offers the red-shirt/colored-shirt case while
asserting the stronger representation-impossibility claim. The active
account retains its valid warning about erasing meaning but not that
unrestricted impossibility. The distinction is between a representation
that preserves the intended lexical relationship and an unrestricted
structural test that treats the two predicates independently.
Adding a genuine meaning constraint differs from adding whatever
conditional would merely force the desired conclusion.

\section{Hamblin and the revival of fallacy theory}

Woods describes Hamblin's 1970 challenge as both historical and theoretical.
Hamblin criticized recycled, implausible textbook examples, traced the
development of fallacy analysis from antiquity through later traditions,
and sought a substantial account of mistaking bad reasoning for good.
His work includes medieval and Indian contributions and a formal
dialectical model. Thus the chapter distinguishes Toulmin's departure
from formal-logical mapping from Hamblin's continued respect for formal
logic and demand that logic recover the fallacies project.

Hamblin's proposed thick theory is not identified with every formal
dialogue result. The chapter recognizes the interest of his dialectical
work without claiming that it delivered the entire theory he called for.
Its comments about the professional literature's response are Woods's
historical assessment at the time of this account, not a present-day
bibliometric census. Notes 9--10 distinguish the foundations-of-arithmetic
project from the possibility of decidability results in particular
formal dialogue and information systems.
In particular, his claim that no fallacy article had appeared in the
\emph{Journal of Symbolic Logic} is not accompanied by a corpus,
search procedure or coverage audit; it is not a validated journal census here.

The Woods--Walton approach, begun in the early 1970s, sought a suitable
existing logic for each fallacy rather than declaring one treatment
canonical. The source identifies uses of intuitionistic, modal,
relatedness, causal, epistemic, interrogative, deontic and dialectical
resources. Its examples are cumulative inquiry intended to block
circularity, aggregate theory applied to composition, and dialogue
games including Circle and DD games. Note 11 explains cumulative inquiry
as precluding retraction of previously supplied answers; note 12 names
the Woods--Walton fragment as another game resource. These are identified
applications, not full game constructions or proofs imported from
the cited collection.

Woods records an objection that the approach gives each fallacy its own
logic and accepts that pluralistic consequence. He also records the
opposite concern that adapting successful preexisting theories can
amount to overconfidence in their appropriateness. Context, purposes
and the inferential role of the parts matter: a familiar surface pattern
is not inherently fallacious simply because of its syntax. The later
paths diverge, with Walton developing schemes/dialogical accounts and
Woods favoring naturalized inference and a causal-response reliabilism
\cite[pp.~370--371 and notes~11--16]{Woods2021ReasoningArgumentation}.

\section{Informal logic: common dissatisfaction, different standards}

The chapter's textbook history includes Copi, Cohen and Nagel, Kahane's
1971 response to demands for relevance, and Johnson and Blair's 1977
work. It connects the movement's institutional development at Windsor,
its journal and conferences with a shift toward the actual contexts of
reason-giving. Woods also warns against treating the earlier textbooks
as wholly worthless and notes that later critics continued to mine them
for motivating examples.

The account does not describe informal logic as one doctrine. Premise
acceptability, premise--conclusion relevance and sufficiency are shared
headings whose details remain disputed. Some approaches emphasize
epistemic justification, others agreement or conflict resolution;
some seek descriptive adequacy, others normative authority. The source's
notes distinguish justified-true-belief and causal-response accounts,
and its cited agenda-relevance and sufficiency work have their own
assumptions. None follows solely from the three headings.
The source also describes resistance to treating first-order mathematical
logic as the best account of actual inference, and to taking a
Tarski-style formal truth semantics as a prototype for every theory of
ordinary argument. An objection to that transfer is not a refutation
of the formal calculus or of Tarski's specified mathematical construction.

Further issues include whether argument and inference coincide, whether
argumentation is essentially attack-and-defense, what empirical sensitivity
a theory needs, whether training transfers to critical thinking, and what
metadialogue or meta-argument adds. Woods's historical and philosophical
survey does not supply a training trial, comparative effect estimate
or uncertainty interval answering those questions. His institutional
contrast between Windsor and Amsterdam is likewise not a measured
claim about every practitioner's activity
\cite[pp.~371--372 and notes~17--22]{Woods2021ReasoningArgumentation}.

\section{The pragma-dialectical model and its critical questions}

Woods presents pragma-dialectics as a model of critical discussion,
contrasting its canonical identity with the doctrinal plurality of
informal logic. He distinguishes its ideal normative model from strategic
maneuvering, which draws on rhetorical considerations to make discussion
persuasive. In that model, \emph{fallacy} means a violation of the specified
discussion rules, a different criterion from every other use of the term.
His history relates its beginning to a simplified counterpart of Barth
and Krabbe's formal dialogue account. His criterion of a canonical theory
combines standing in the marketplace of ideas with a durable approving
readership; neither is by itself proof of correctness or universal coverage.

The chapter advertises ten critical-discussion rules but prints the
following content in nine grouped bullets. These are the actual groups
in Woods's presentation, not an independently reconstructed ten-rule
original with its numbering silently changed.
\begin{enumerate}
\item Parties may advance and challenge standpoints.
\item A party advancing a standpoint bears the corresponding proof burden.
\item Attacks and defenses address the relevant topic.
\item Parties do not disown their attributed implicit premises or
      misattribute implicit premises to the other party.
\item Parties respect the accepted starting premises rather than
      introducing an unaccepted premise or denying an accepted one.
\item A conclusive defense uses an appropriate scheme correctly.
\item A conclusive argument is valid, or becomes so by making its
      implicit premises explicit.
\item An unsuccessful defense is withdrawn; a successful defense calls
      for withdrawal of the challenge.
\item Expression is clear, and interpretation of the other party is fair.
\end{enumerate}
The source itself asks whether the rule set satisfies its own clarity
requirement. Its reference to \emph{rule 9} is retained as a reference in
that discussion, not used to identify a rule in another numbering
\cite[p.~372]{Woods2021ReasoningArgumentation}.

Woods's criticisms raise seven connected questions: whether the rules
define critical discussion by necessary and sufficient conditions;
whether actual discussions instantiate that model; what follows from
empirical inadequacy; whether idealization can rescue it; what establishes
the ideal model's normative authority; whether the deductivist condition
permits arbitrary validation; and how well the model transfers. His
negative transfer judgment mentions collective bargaining or treaty
negotiation, parliamentary and political debate, health-policy discussion
and visual argumentation. The chapter gives no comparative task samples,
measured outcomes or uncertainty for that judgment. It is an assessment
by the author, not an empirical transfer result established here.

\begin{hbtheorem}{thm:p18-r56-conditional-repair}{proved=annex:p18.r56.conditional-repair}{Arbitrary deductivization adds the disputed support}
Let \(\Gamma=\{\gamma_1,\ldots,\gamma_n\}\) be any finite set
of classical propositional premises, allowing \(n=0\), and let
\(\varphi\) be a conclusion. Put
\[
 A_\Gamma=\bigwedge_{i=1}^n\gamma_i,\qquad
 \chi_\Gamma=A_\Gamma\to\varphi,
\]
where the empty conjunction is true. Then
\(\Gamma\cup\{\chi_\Gamma\}\models\varphi\).
If the original argument is invalid, every valuation witnessing
\(\Gamma\not\models\varphi\) makes \(\chi_\Gamma\) false.
Thus the added premise's acceptability is not established by the
formal validity it creates. An independently warranted implicit
premise is a different case.
\end{hbtheorem}

The source does not merely dismiss the underlying dialogical idea.
It suggests that its appeal can reflect genuine features of giving
reasons to someone who disagrees, even if the particular theory handles
them inadequately. A defensible idealization must specify the relation
between model and practice, the omitted features, and the standard by
which it claims authority; those are not supplied by the word \emph{ideal}.

\section{Schemes, law, psychology, and computational representations}

The chapter's newer developments include legal reasoning, logical
information dynamics and argument networks, rapprochement with psychology,
and scheme-based software/AI work. A scheme is a pattern of reason-giving
with evaluation questions, including accuracy of premises, unexpressed
load-bearing assumptions and allowable exceptions. The source additionally
asks under which assumption the pattern would succeed and whether that
assumption may reasonably be treated as provisionally true.

Presumption and defeasibility must not automatically be identified.
Woods notes Walton's tendency to relate them, but also a contrasting
view: defeasibility concerns openness to contrary information, whereas
presumption can describe a disposition to regard defeating circumstances
as surprising. The source reports 29 schemes in Walton's 1996 presentation
and 96 in the 2008 successor. These are reported source-specific counts,
not proof of an exhaustive taxonomy or a count of completed handbook
treatments. Storing a finite list computationally does not by itself
solve interpretation, classification or evidential assessment.

Computational modeling requires explicit representations of the proposed
account and suitable input data. The source describes adding answers to
scheme questions as premises in a formal model. That operation must not
collapse into the arbitrary repair of the preceding theorem. A negative
answer, unresolved exception or unsupported assumption may instead
require narrowing, suspending or withdrawing the conclusion. The question
of whether the model represents the phenomenon, or the data have been
altered to fit the model, remains substantive.

Woods closes with questions about where schemes originate, what gives
them probative force, what motivates their critical questions, and how
one knows that a question list is complete. Identifying a recognizable
pattern does not answer those questions
\cite[pp.~373--374 and notes~23--24]{Woods2021ReasoningArgumentation}.

\section{Comparison and a bounded assessment method}

The existing scheme reconstruction method marks a premise as an
attributed implicit commitment only when text or context supports it;
otherwise it calls it needed but unestablished. It also forbids rescuing
an argument merely by asserting that its premises imply its conclusion
(see \cref{def:p18-scheme-reconstruction-method}).
The separate purpose, material relation and reasoning-form dimensions
at \cref{def:p18-scheme-three-dimensions} help express Woods's concern
that a selected manageable class should not become a universal theory.
An explanatory hypothesis, an action recommendation and a strict
deduction can use similar subject matter but require different standards.

The dialogue-purpose treatment distinguishes inquiry, persuasion,
negotiation, deliberation and other episodes, including changes of
purpose within an exchange (see \cref{def:p18-dialogue-purpose}).
It therefore does not presume that every disagreement is a critical
discussion with one idealized closing criterion. An acceptable move
in a negotiation need not be evidence that a proposition is true.

Hahn and Collins's 5.5 distinction between objects/processes and
procedural/epistemic norms overlaps this account, but is not identical
to having/spotting/drawing or as-such/in-use
(see \cref{def:p18-r55-objects-processes,def:p18-r55-norms}).
Their probabilistic perspective supplies a stated measure of evidential
change under a model; Woods asks which agent-centered process and standard
that model represents. A formal conditional probability, a valid
entailment, a dialogical obligation and a useful actual inference
cannot substitute for one another without a justified bridge.
His challenge to uninstantiable ideals does not refute the explicitly
conditional finite calculations or task-bound empirical records in 5.5.

\begin{hbmethod}{def:p18-r56-use-method}{Select the question before selecting its standard}
Record the actual episode and intended conclusion: a truth claim,
probability judgment, explanation, recommendation or resolution of a
disagreement. Separate the static premise--conclusion relation from
recognition, belief revision and the action of drawing a conclusion.
State the model class, background meaning and ordinary facts that a
formalization preserves; do not erase a lexical connection and then
mistake its lost consequence for a defect of the original.

Identify the scheme and dialogue rules only for their declared scope.
For an implicit premise, give grounds for attribution and for its
acceptability; test whether it merely restates the desired implication.
For a defeasible conclusion retain exceptions and uncertainty rather
than presenting a formalized default as an unconditional fact.
For a claimed cognitive or training result, require its actual task,
sample, design, response measure and uncertainty; a conceptual example
or historical judgment supplies no substitute.

Finally distinguish a philosophical proposal, a formal result, a
source-attributed historical assessment and an empirical finding.
Report the argument's supported force and what remains unresolved.
The procedure is an editorial synthesis of the stated distinctions,
not a complete decision algorithm or an empirical accuracy guarantee.
\end{hbmethod}
